\subsection{THE EXTENDED LINE}

We shall now define, by adjoining two new elements to R, a topological space \(\overline{R}\) such that every homeomorphism of R onto a bounded open interval I of R can be extended to a homeomorphism of \(\overline{R}\) onto the closed interval having the same end-points as I.

Let \(\overline{\mathbf{R}}\) be the set obtained by adjoining (Set Theory, R, § 4, no. 5) two new elements to \(\mathbf{R}\) , denoted by \(-\infty\) and \(+\infty\) . We extend the ordering on \(\mathbf{R}\) to \(\overline{\mathbf{R}}\) by putting \(-\infty < a\) and \(a < +\infty\) for all \(a \in \mathbf{R}\) , and \(-\infty < +\infty\) ; it is clear that we thus obtain a linearly ordered set, whose ordering induces the ordering of the real line on \(\mathbf{R}\) . Next, consider the topology on \(\overline{\mathbf{R}}\) generated by the set of open intervals of \(\overline{\mathbf{R}}\) . Since the trace on \(\mathbf{R}\) of an open interval of \(\overline{\mathbf{R}}\) is an open interval of \(\mathbf{R}\) , this topology induces on \(\mathbf{R}\) the topology of the real line.

DEFINITION 1. The set \(\overline{\mathbf{R}}\) endowed with the order structure and the topology defined above is called the extended real line.

When using the extended line \(\overline{\mathbf{R}}\) it is often convenient to refer to its points as real numbers, by abuse of language; the points of \(\mathbf{R}\) are then called finite real numbers. We shall adopt this convention in this section and the three following sections of this chapter; whenever we adopt this convention in future we shall indicate explicitly to which part of the text it extends.

If \(a\) is a finite real number, the intervals \([a, +\infty[ \text{ and } ] - \infty, a]\) (resp. \(]a, +\infty[ \text{ and } ] - \infty, a[\) ) of \(\overline{\mathbf{R}}\) are contained in \(\mathbf{R}\) and coincide with the intervals of \(\mathbf{R}\) hitherto denoted by \([a, \rightarrow[ \text{ and } ] \leftarrow, a]\) (resp. \(]a, \rightarrow[ \text{ and } ] \leftarrow, a[\) ); this new notation is much more often used. Again, \(\mathbf{R}\) coincides with the interval \(] - \infty, +\infty[ \text{ of } \overline{\mathbf{R}}, \text{ and is sometimes so denoted.}\)

PROPOSITION 2. Every homeomorphism \(f\) of \(\mathbf{R}\) onto an interval \(]a, b[\) can be extended to a homeomorphism \(\overline{f}\) of \(\overline{\mathbf{R}}\) onto \([a, b]\) . If \(f\) is an increasing function, then \(\overline{f}\) is an order isomorphism of \(\overline{\mathbf{R}}\) onto \([a, b]\) .

Let \(f\) be an increasing homeomorphism. If we extend \(f\) to \(\mathbf{R}\) by putting \(\overline{f}(-\infty) = a\) and \(\overline{f}(+\infty) = b\) , it is obvious that \(\overline{f}\) is a strictly increasing mapping (and therefore a bijection) of \(\overline{\mathbf{R}}\) onto \([a, b]\) . Hence \(\overline{f}\) maps every open interval of \(\overline{\mathbf{R}}\) onto an interval which is open with respect to \([a, b]\) , and is therefore a homeomorphism of \(\overline{\mathbf{R}}\) onto \([a, b]\) by reason of Definition 1 and Proposition 5 of § 1, no. 4.

If \(f\) is decreasing, we apply what has been proved to the increasing homeomorphism \(x \to -f(x)\) of \(\mathbf{R}\) onto \(] - b, - a[\) .

All the properties of the interval \([a, b]\) obtained in §2, which involve only the order structure and the topology of the interval, can therefore be transported to \(\overline{R}\) ; hence the following propositions:

PROPOSITION 3. The extended real line is compact.

Hence (Chapter II, § 4, no. 1, Theorem 1) there is a unique uniformity on \(\overline{\mathbf{R}}\) compatible with its topology; this uniformity is isomorphic with the uniformity induced on \([a, b]\) by the additive uniformity of \(\mathbf{R}\) . But it should be remarked that the uniformity induced on \(\mathbf{R}\) by that of \(\overline{\mathbf{R}}\) is not the additive uniformity of \(\mathbf{R}\) (although it is compatible with the topology of the real line); for \(\mathbf{R}\) is a complete space with respect to its additive uniformity, but is not a complete subspace of \(\overline{\mathbf{R}}\) , since it is not closed in \(\overline{\mathbf{R}}\) .

PROPOSITION 4. Every non-empty subset of \(\overline{\mathbf{R}}\) has a least upper bound and a greatest lower bound.

The least upper bound (resp. greatest lower bound) of a non-empty subset A of R is denoted by sup A (resp. inf A). Clearly

\[
\inf A \leqslant \sup A.\tag{1}
\]

If \(\mathbf{A} \subset \mathbf{B}\) , then \(\sup \mathbf{A} \leqslant \sup \mathbf{B}\) and \(\inf \mathbf{A} \geqslant \inf \mathbf{B}\) (Set Theory, Chapter III, § 1, no. 9, Proposition 4).

PROPOSITION 5. A subset A of \(\overline{R}\) is connected if and only if A is an interval.

COROLLARY. The extended real line is a connected, locally connected space.

PROPOSITION 6. A mapping \(f\) of an interval \(I\) of \(\overline{R}\) into \(\overline{R}\) is a homeomorphism of \(I\) onto \(f(I)\) if and only if \(f\) is strictly monotonic and continuous on \(I\) ; \(f(I)\) is then an interval of \(\overline{R}\) .

Finally, the functions \(\sup (x,y)\) and \(\inf (x,y)\) are continuous on \(\overline{\mathbf{R}}\times \overline{\mathbf{R}}\)

